Predicting how much a trained neural network has overfit --- its generalization gap --- is more tractable
from weight tensors than predicting its test accuracy outright. Motivated by this counterintuitive
asymmetry, we conduct the first controlled comparison of equivariant and non-equivariant weight-space
encoders on generalization gap as the primary prediction target, using a dual-target design on the
Unterthiner CIFAR-10 CNN zoo. We find that gap is learnable at Spearman $r > 0.5$ with standard encoders,
and that architecture specificity matters: only Neural Functional Transformers, with their cross-layer
attention, improve over the flat baseline on gap ($r = 0.575$), while within-layer and graph-structured
equivariant architectures underperform it. Most strikingly, a partial Spearman analysis reveals that
NFT gap predictions carry substantial information about true gap that is independent of test accuracy ---
confirming that generalization gap and test accuracy encode distinct signals in weight space. We report
both the positive findings and a null result on the target-specificity mechanism hypothesis, with
identified confounders and proposed corrective experiments, providing a transparent empirical foundation
for gap-targeted weight-space learning.
